% reality/data-DE.tex

\chapter{Messdaten und Beobachtungskontext}
\label{chap:reality-data}

\index{Beobachtung!Datenquellenkontext}
\index{Datenquelle}

Messdaten bilden die Beobachtungsebene von AIM. Sie liefern Evidenz über
physische Infrastruktur und Laufzeitzustand, definieren aber für sich allein
keine Engineering Identity.

\section{Rolle der Daten}

\begin{definition}[Datenquelle]
\index{Datenquelle!Definition}
Eine Datenquelle ist jeder Ursprung geometrie- oder zustandsbezogener
Beobachtungen, die in der Eisenbahnmodellierung verwendet werden.
\end{definition}

Typische Quellen sind geodätische Vermessungen, Gleisgeometriemessungen,
historische ingenieurtechnische Datensätze, strukturierte Austauschformate und
aus Punktwolken abgeleitete Produkte.

\begin{principle}[Beobachtungsprinzip]
\index{Beobachtungsprinzip}
Messdaten sind ingenieurtechnische Beobachtungsevidenz und müssen anhand
ausdrücklicher Modell- und Kontextannahmen interpretiert werden.
\end{principle}

\section{Dateneigenschaften}

Reale Eisenbahndaten sind heterogen und weisen typischerweise
Unvollständigkeit, Rauschen, Mehrdeutigkeit und quellübergreifende
Inkonsistenz auf. Daten können deshalb nicht unmittelbar mit konstruktiver
Identität, metrischer Realisierung oder realisiertem Zustand gleichgesetzt
werden.

\begin{definition}[Messunsicherheit]
\index{Unsicherheit!Messung}
Messunsicherheit beschreibt die Abweichung zwischen beobachteten Größen und
den entsprechenden physischen Größen.
\end{definition}

Unsicherheit pflanzt sich durch Interpretation, Rekonstruktion und
Laufzeitbewertung fort.

\section{Interpretationsebene}

\begin{definition}[Interpretation]
\index{Interpretation!Beobachtungsdaten}
Interpretation ist der Prozess, beobachteten Daten in einem bestimmten Kontext
geometrische und semantische Bedeutung zuzuweisen.
\end{definition}

Interpretation umfasst Segmentierung, Topologiezuordnung,
Stationierungskontext und die Verknüpfung mit Längsprofil-, Überhöhungs- und
Laufzeitgrößen. Der Prozess ist kontextabhängig und im Allgemeinen nicht
eindeutig.

\section{Von Daten zu strukturierten Modelleingaben}

Die Datenverarbeitung in AIM bildet ab:
\[
\text{Beobachtungsdaten}
\rightarrow
\text{interpretierte Modelleingabe}
\rightarrow
\text{Kontext der Laufzeitbewertung}.
\]

\begin{definition}[Rekonstruktion aus Beobachtungsdaten]
\index{Rekonstruktion!aus Beobachtungsdaten}
Rekonstruktion ist der Prozess, aus beobachteten Daten strukturierte
Modelleingaben abzuleiten und dabei Provenienz und Unsicherheitskontext zu
bewahren.
\end{definition}
\index{Provenienz!Beobachtungsrekonstruktion}

Diese Rekonstruktion kann Filterung, Segmentierung, Ausrichtung an
Stationierung und Topologie sowie eine mit dünn besetzten Modellen kompatible
Approximation umfassen.

\section{Geometrische und topologische Evidenz}

Beobachtungsdatensätze können geometrische Evidenz -- Positionen, Richtungen
und krümmungsbezogene Stichproben -- sowie topologische Evidenz --
Konnektivität, Verzweigungs- und Routenkontext -- enthalten. Diese Ebenen sind
begrifflich getrennt zu halten und zugleich gemeinsam zu interpretieren.

\section{Container- und Austauschkontext}

Containerformate können gemischte Geometrie, Metadaten und semantische Inhalte
tragen. AIM behandelt Containerinhalte daher als Import-Evidenz, die vor
Laufzeitbewertung oder Realisierungsvergleich in ausdrückliche Strukturen der
Modellebenen normalisiert werden muss.

\section{Verbindung zu AIM}

\begin{summary}
Im Teil Realität werden Daten als Beobachtungskontext für
realisierungsbewusste und laufzeitbewusste Interpretation behandelt.

Sie sind weder konstruktive Identität noch abschließende ingenieurtechnische
Wahrheit. Ingenieurtechnische Interpretation entsteht erst nach
ausdrücklicher Rekonstruktion, Kontextzuordnung und dem Vergleich von Ziel und
Realisiert.
\end{summary}
